
\subsection{Preregistered zero-delay numerical check}
Five Julia/PowerDynamics integrations comprise one independent replay of the
legacy 1 MW event and four physical-variant events. The fixed candidate has
94.215489\% GFL and 312.523288 MW retained SG;
no design variable is retuned. The trim residual is
$6.574\times10^{-12}$, and all 144 finite poles,
including one gauge, match bijectively with numerical error
$0$. The physical rightmost pole is
$-0.050057324\,\mathrm{s}^{-1}$.

\begin{center}\begin{tabular}{lrrl}\toprule
ZIP event & $\max|\Delta f|$ [Hz] & RoCoF [Hz/s] & Both limits\\\midrule
bus16\_1MW & 0.005000 & 0.001670 & pass \\
bus16\_100MW & 0.468755 & 0.171772 & pass \\
bus8\_100MW & 0.410549 & 0.512750 & fail \\
bus29\_100MW & 0.588571 & 0.281888 & fail \\
\bottomrule\end{tabular}\end{center}
These are sampled maxima at generator buses 30--39, measured by a causal
0.5 s phase window over 0--61 s, with limits 0.5 Hz and 0.5 Hz/s.
The event is a nominal ZIP-setpoint increment, not constant realized demand.
The two failed 100 MW events remain failed after correcting the energy port.

The positive-energy gradient-times-RHS identity has maximum residual
$2.919\times10^{-10}$ MW. Independently integrating stored
traces at 5 ms leaves maximum energy and frequency-area residuals of
0.046437 MJ and
0.046402 MJ; these are exposed quadrature
errors, not interval certificates. The SG35 governor is directly observed
on its limiting branch in 380 samples of the
bus29 event, giving approximately 1.900 s
of sampled limiting and an antiwindup integral of
-0.728187 MJ. Ignoring it changes the area balance.
The isolated campaign is reproducible from
\texttt{experiments/physical\_frequency\_bridge\_20261004/}.
All frozen source and candidate hashes were checked unchanged.
This closes the physical-port audit for the new variant. It does not close
the full-model replacement-floor or global-optimality gates.
